% ============================================================================
%  Chapter 6. Architecture, and the Falsification of Three Claimed Contributions
%
%  \input file. Requires booktabs, amsmath, natbib in the parent preamble.
%  Bibliography: accord/thesis/refs.bib
%
%  Traceability: accord/CORRECTIONS.md C1, C2, C4, C5, C7, C8;
%  accord_phase3_v2_ablations.py; transfer_results_correct.json;
%  the Stage 0 pretraining logs.
% ============================================================================

\chapter{Architecture, and the Falsification of Three Claimed Contributions}
\label{ch:arch}

\section{What This Chapter Does}
\label{sec:arch-overview}

The preceding chapters concerned data and supervision. This chapter concerns the
model, and it is written in an unusual form: as a falsification report. ACCORD
was specified with four purpose-built components layered over a biomedical
encoder, and three claimed contributions were attached to them. An audit of the
code that produced the reported results establishes that none of the three
survives.

\begin{enumerate}
\item \textbf{Claimed:} the model discovers clinical note priority without
      supervision, converging to weights $\alpha_N \approx 0.5$,
      $\alpha_D \approx 1.0$. \textbf{Actual:} the parameter is
      \emph{initialised} at exactly $[0.5, 1.0]$ and moves by $1.8\%$ and
      $0.6\%$ across seeds. Section~\ref{sec:c1}.
\item \textbf{Claimed:} an auxiliary rule-type head reaches $F_1 = 0.9603$,
      evidencing pharmacologically grounded structure. \textbf{Actual:} the rule
      identifier is written into the model's input features and predicted in the
      same forward pass, one linear layer away. Section~\ref{sec:c2}.
\item \textbf{Claimed:} guideline-grounded harm estimation produces ranked
      danger scores. \textbf{Actual:} the guideline bank is twenty random
      vectors, and the harm head receives no loss term at any point in training.
      Section~\ref{sec:c3}.
\end{enumerate}

A fourth finding, the collapse of self-supervised pretraining to the
constant-prediction solution, is reported in Section~\ref{sec:stage0} together
with a retraction of the explanation originally given for it.

The chapter closes with the constructive residue: a set of diagnostics, each
costing a few lines of code, any one of which would have caught the
corresponding failure before it reached a results table.

\section{Backbone and Adaptation}
\label{sec:backbone}

The encoder is BioMedBERT \citep{gu2021biomedbert}, a $110$M-parameter
transformer pretrained from scratch on PubMed abstracts and full texts. It was
selected over ClinicalBERT \citep{alsentzer2019clinicalbert} on the basis of
published benchmark performance; that choice is defensible and is not among the
findings of this chapter.

Adaptation uses LoRA \citep{hu2022lora} with rank $r = 16$ and
$\alpha = 32$ on the attention projections. Chapter~\ref{ch:eval} reports the
one place in this project where the backbone and the adapter were evaluated
correctly, on MedNLI \citep{romanov2018mednli}, and the result there is sound.
It is also, as that chapter argues at length, a result about BioMedBERT, not
about ACCORD.

Input construction: each candidate's nursing sentence and discharge sentence are
encoded, each truncated to $256$ WordPiece tokens taken from the head of the
sequence. Section~\ref{sec:stage0} returns to that truncation.

\section{The Specified Architecture}
\label{sec:spec}

\begin{table}[t]
\centering
\small
\begin{tabular}{llr}
\toprule
\textbf{Component} & \textbf{Stated function} & \textbf{Params} \\
\midrule
CDE  & [TODO: expand acronym from the original specification] & 527,104 \\
ATCG & [TODO: expand acronym from the original specification] & 395,269 \\
PPC  & Pharmacological Prior Conditioning: injects drug-class, & 2,496 \\
     & interaction, and contraindication priors & \\
GGHE & Guideline-Grounded Harm Estimation: retrieval-augmented & 614,916 \\
     & harm scoring against published HF guidelines & \\
\midrule
     & \textbf{Total new parameters (v1 specification)} & \textbf{1,539,785} \\
\bottomrule
\end{tabular}
\caption{The specified architecture, as described in the project's own
documentation. Two acronym expansions are not recoverable from surviving
artifacts and are marked [TODO] rather than reconstructed. The total,
$\approx 1.54$M, is the figure quoted throughout the project's writeups.
Section~\ref{sec:actual} establishes that this table describes v1 and that no
reported result was produced by v1.}
\label{tab:spec}
\end{table}

Table~\ref{tab:spec} is what the project claimed to have built. The parameter
counts are internally consistent and sum to the quoted $1.54$M. The problem is
not arithmetic.

\section{The Architecture That Actually Ran}
\label{sec:actual}

\begin{table}[t]
\centering
\small
\begin{tabular}{p{0.24\linewidth}p{0.68\linewidth}}
\toprule
\textbf{Claimed} & \textbf{Actual, in the code that produced the results} \\
\midrule
PPC, $2{,}496$ params, injects pharmacological priors
& \textbf{Deleted in v2.} Replaced by a function \texttt{\_prior()} returning a
  $64$-dimensional vector of zeros with two positions overwritten:
  \texttt{pf[i,0]} $=$ a danger multiplier and \texttt{pf[i,1]} $=$
  \texttt{rule/4.0}. \textbf{Zero learned parameters.} No drug class, no
  interaction table, no contraindication list. \\
\addlinespace
GGHE, retrieval-augmented scoring against published HF guidelines
& A single \texttt{nn.Parameter(torch.randn(20, 256))}: twenty random vectors,
  cross-attended from one query. No corpus, no index, no retrieval, no guideline
  text. The code comment reads \emph{``In full version this attends over
  FAISS-retrieved chunks.''} \\
\addlinespace
Guideline-grounded harm estimation produces ranked danger scores
& \texttt{harm\_head} and \texttt{explanation\_head} \textbf{receive no loss
  term}. The total objective is
  $\mathcal{L} = 1.0\,\mathcal{L}_{\text{contra}} + 0.3\,\mathcal{L}_{\text{rule}}$.
  Neither head appears in it. \\
\addlinespace
$1.54$M new parameters
& Describes v1. v2 adds an \texttt{AsymmetricCrossNoteEncoder}: two $8$-head,
  $768$-dimensional multi-head attention blocks plus two $768 \times 768$
  projections, $\approx 5.9$M parameters, and removes PPC's $2{,}496$. The model
  that produced every reported result carries $\approx 7.5$M new parameters, not
  $1.54$M. \\
\bottomrule
\end{tabular}
\caption{Specification against implementation. The right column is what the
code does. Every result reported in this thesis, and every result reported in
the project's earlier writeups, was produced by the right column.}
\label{tab:actual}
\end{table}

Table~\ref{tab:actual} should be read as a whole rather than component by
component, because the pattern matters more than any single row. In each case
a component was specified, a placeholder was written so that the surrounding
code would run, and the placeholder was never replaced. The pipeline executed
end to end, produced numbers, and those numbers were reported against the
specification rather than against the placeholder. Nothing in the codebase
compared the two.

This is the structural failure of the project. It is a documented failure mode
of machine learning systems \citep{sculley2015debt}, and the remedy is
mechanical: a model summary printed at training time, listing every module, its
parameter count, and whether it appears in the loss graph.
Section~\ref{sec:diagnostics} states it as a protocol.

\section{Falsification 1: Authority Conditioning Is an Initialisation Artifact}
\label{sec:c1}

\subsection{The claim}

The project's third stated contribution was that ACCORD recovers clinical note
priority without ever being told it. Specifically: without explicit authority
labels, the model's authority weights converge to the correct clinical
hierarchy, $\alpha_N \approx 0.5$ for nursing and $\alpha_D \approx 1.0$ for
discharge. This was presented as evidence that institutional note priority is
learnable from weak supervision alone.

\subsection{The code}

\begin{quote}\ttfamily\small
def \_\_init\_\_(self, hidden\_size=768, num\_heads=8, dropout=0.1,
authority\_dampening=0.5): \\
\hphantom{def }self.authority\_scale = nn.Parameter(torch.tensor([authority\_dampening, 1.0]))
\end{quote}

The parameter is initialised at exactly $[0.5, 1.0]$: the values reported as
discovered.

\subsection{The measurement}

Across three seeds the parameter moved to $(0.509, 1.006)$, relative changes of
$1.8\%$ and $0.6\%$. Third-decimal agreement across independent seeds is not the
signature of a recovered signal. It is the signature of a parameter that did not
train.

\subsection{The mechanism}

The scale is applied to a cross-attention branch that is immediately
LayerNorm'd. LayerNorm is invariant to a per-branch scalar rescaling of its
input up to the effect on its own learned affine parameters, so the normalisation
removes most of what the authority scale does, and the gradient flowing back to
it is correspondingly near zero. The parameter is not merely uninformative; it
is architecturally prevented from mattering.

\subsection{The project's own ablation already refuted the claim}

The clearest evidence was present in the results table and was not read. A
symmetric variant, with $\alpha_N = \alpha_D$, \emph{outperforms} the full
asymmetric model:

\begin{table}[h]
\centering
\small
\begin{tabular}{lrr}
\toprule
 & \textbf{Contradiction $F_1$} & \textbf{Accuracy} \\
\midrule
Symmetric (constrained $\alpha_N = \alpha_D$) & \textbf{0.7775} & \textbf{0.9192} \\
Full asymmetric model                          & 0.7368 & 0.8939 \\
\bottomrule
\end{tabular}
\caption{The asymmetry ablation. Removing the mechanism claimed as a
contribution improves both reported metrics. Both rows are measured against
\emph{weak} labels on the held-out split, not against pharmacist labels, so they
describe agreement with an $84.9\%$-positive teacher rather than clinical
correctness. [TODO: report the majority-class baseline on this exact split.
Without it, an $F_1$ of $0.7775$ against supervision that is $84.9\%$ positive
cannot be interpreted, because a constant-positive predictor achieves a high
$F_1$ by construction under that marginal.]}
\label{tab:acne}
\end{table}

\subsection{What should be reported instead}

The claim is retracted in full. Its replacement is the diagnostic that would
have caught it, stated as a general rule: \emph{for any parameter reported as
learned, publish $\lVert\theta_{\text{final}} - \theta_{\text{init}}\rVert /
\lVert\theta_{\text{init}}\rVert$ and a gradient-norm trace.} Applied here,
those two quantities are $1.8\%$/$0.6\%$ and approximately zero, and the claim
does not survive its own first diagnostic.

\subsection{What this does not establish}

It does not establish that clinical note priority is unlearnable, nor that
asymmetric conditioning is a bad idea. It establishes that this implementation
did not test the hypothesis. The hypothesis remains open, and a correct test
would require a parameter placed outside a normalisation that erases it, an
initialisation that does not encode the expected answer, and a measurement of
displacement from that initialisation.

\section{Falsification 2: The Rule-Type Head Is Label Leakage}
\label{sec:c2}

\subsection{The claim}

A multi-task auxiliary head predicts which of the four rules generated a
candidate, reaching $F_1 = 0.9603$ (reported as $0.953 \pm 0.006$ in the
project's ablation table). This was offered as evidence that the model learns
pharmacologically grounded structure rather than surface cues.

\subsection{The code}

The prior-feature constructor writes the rule identifier into the model's input:

\begin{quote}\ttfamily\small
pf[i, 1] = rule / 4.0\hfill\# rule comes from the batch \\
\ldots \\
return self.gghe(edge\_repr, prior) \\
rule\_logits = self.rule\_type\_head(attended)\hfill\# predicts RULE\_MAP[rule]
\end{quote}

The auxiliary label is an input feature and the prediction target in the same
forward pass, separated by one linear layer. The task is a copy operation, and
$F_1 = 0.9603$ is the accuracy of that copy.

\subsection{The independent confirmation}

Leakage of this kind is invisible in domain and unmistakable out of domain,
because the leaked feature is present in both cases but the surface correlations
supporting a genuine solution are not. Table~\ref{tab:transfer} reports the
transfer test.

\begin{table}[t]
\centering
\small
\begin{tabular}{lr}
\toprule
\textbf{Setting} & \textbf{Rule-type $F_1$} \\
\midrule
Heart failure (in domain) & 0.9603 \\
Sepsis (transfer)         & 0.0654 \\
AKI (transfer)            & 0.0595 \\
\midrule
Rule-prediction accuracy on the held-out pharmacist set & 45.2\% \\
\bottomrule
\end{tabular}
\caption{Rule-type prediction in domain and under transfer. A head reaching
$0.96$ in domain collapses to $0.065$ and $0.060$ on two other ICD-9 cohorts
drawn from the same institution, the same note types, the same extractor, and
the same four rules. A drop of that magnitude between populations that differ
only in index diagnosis is not distribution shift in any ordinary sense. It is
the removal of whatever the head was actually using.}
\label{tab:transfer}
\end{table}

The transfer collapse is the better number to report, and it is reported here in
place of $0.9603$, which is withdrawn. It is the \emph{evidence} of the leakage
rather than an embarrassment to be minimised, and it happens to answer the
generalisation question directly: this model does not transfer between clinical
cohorts.

\subsection{Positioning}

Auxiliary-label leakage of this shape belongs to a well-documented family. NLI
models exploit hypothesis-only artifacts \citep{poliak2018hypothesis,
gururangan2018artifacts} and syntactic heuristics \citep{mccoy2019hans};
radiology models key on hospital-identifying image features and collapse across
sites \citep{zech2018generalization}; the general phenomenon is shortcut
learning \citep{geirhos2020shortcut}. The contribution of this section is not
the observation that shortcuts exist. It is the specific, checkable point that
\emph{an auxiliary head whose target is derivable from the input features is not
an auxiliary task}, and that the cheapest test for it is out-of-domain transfer,
not held-out accuracy.

\section{Falsification 3: Harm Estimation Was Never Trained}
\label{sec:c3}

\subsection{The claim}

The project posed a research question about guideline-grounded severity ranking:
whether contradictions can be ordered by clinical danger using evidence from
published heart-failure guidelines. GGHE was the component that would answer it.

\subsection{The code}

Three facts, each independently sufficient.

\paragraph{There is no guideline corpus.} The ``guideline bank'' is
\texttt{nn.Parameter(torch.randn(20, 256))}. Twenty vectors, drawn from a
standard normal at initialisation, cross-attended from a single query. There is
no document collection, no index, no retrieval step, and no text of any kind. The
code comment states that a full version would attend over retrieved chunks; that
version does not exist.

\paragraph{The harm head receives no gradient.} The total objective is
$\mathcal{L} = 1.0\,\mathcal{L}_{\text{contra}} +
0.3\,\mathcal{L}_{\text{rule}}$. The harm head and the explanation head are not
in it. They are randomly initialised modules whose outputs are computed,
recorded, and reported, and whose weights never move.

\paragraph{The outputs are consistent with an untrained sigmoid.} Across the
$299$ gold pairs, harm scores span $0.398$ to $0.610$: a narrow band centred on
$0.5$, which is what a randomly initialised head produces. The derived
high-risk rate is $0.0$ in all three cohorts, because no score ever crosses the
threshold.

\subsection{Consequence}

\textbf{The severity-ranking research question was not tested.} The component
that would have answered it was never trained, and the scores it emits are
random-initialisation noise. Chapter~\ref{ch:eval} reports the corresponding
empirical result, a Spearman correlation between harm score and pharmacist
danger rating of $-0.0001$ at $p = 0.999$, which is as close to a pure null as a
finite sample produces. That number is not a negative finding about severity
ranking. It is a measurement of an untrained module, and it should be read as
confirming Section~\ref{sec:c3} rather than as evidence about the underlying
question.

Stating this correctly matters more than it might appear. A reader encountering
$\rho = -0.0001$ without this section would reasonably conclude that clinical
severity is not predictable from these features. The correct conclusion is that
nothing was asked of the model.

\section{The Stage 0 Collapse}
\label{sec:stage0}

\subsection{The setup}

Before supervised training, ACCORD was to be pretrained with a self-supervised
cross-note coherence objective: discriminate genuine (nursing, discharge) pairs
from mismatched pairs, on the theory that a representation able to tell a real
pairing from a fake one has learned something about cross-note consistency.

\subsection{The result}

The objective converged to exactly $\ln 2 = 0.693$ with development accuracy of
$0.5000$ at every epoch. This is the constant-prediction solution: binary
cross-entropy at $p = 0.5$ for every example. The model learned to emit a
constant.

Critically, this held in \emph{both} training configurations:

\begin{table}[h]
\centering
\small
\begin{tabular}{lrrr}
\toprule
\textbf{Configuration} & \textbf{Trainable params} & \textbf{Final loss} & \textbf{Dev acc.} \\
\midrule
Frozen backbone             & 592,647    & 0.693 ($=\ln 2$) & 0.5000 \\
Partially unfrozen backbone & 22,446,855 & 0.693 ($=\ln 2$) & 0.5000 \\
\bottomrule
\end{tabular}
\caption{Stage 0 coherence pretraining. Increasing trainable capacity by a
factor of $37.9$ changes nothing. Identical collapse under wildly different
capacity is strong evidence that the constraint is architectural or
representational, not a matter of insufficient parameters or optimisation
budget.}
\label{tab:stage0}
\end{table}

\subsection{A retracted explanation}

The original diagnosis was that clinical notes within an institution are too
homogeneous for a coherence objective to separate real from mismatched pairs.
\textbf{This explanation is withdrawn.} Nothing in the pipeline measured note
similarity, within an admission or between admissions. No embedding distance, no
lexical overlap statistic, no discriminability estimate was computed. The
explanation was asserted, it was plausible, and it was never tested. It is
withdrawn not because it is known to be false but because there is no evidence
for it, and an untested causal story attached to a null result is the most
easily propagated kind of error.

\subsection{Two mechanisms that are visible in the code}

Two properties of the implementation are directly observable and are each
sufficient to produce the observed collapse. Neither was isolated
experimentally, and the language here is chosen to reflect that.

\paragraph{A one-dimensional bottleneck.} The scoring path is
\texttt{conflict\_scorer} $\rightarrow$ \texttt{sigmoid} $\rightarrow$
\texttt{Linear(1,1)}. The entire pair representation is compressed to a single
scalar, squashed into $(0,1)$, and then passed through an affine map of one
parameter. Three consequences follow. The model can express only a monotone
function of one scalar, so any structure requiring two or more directions is
unrepresentable. The sigmoid sits between the representation and the loss, so
whenever the scalar drifts into a saturated region the gradient reaching
everything upstream is attenuated. And a scalar that collapses to a constant
produces a constant output regardless of how much capacity sits behind it, which
is exactly the invariance observed in Table~\ref{tab:stage0}.

\paragraph{Head truncation into the de-identification header.} Each note is
truncated to its first $256$ WordPiece tokens, taken from the head. As
Chapter~\ref{ch:data} noted, the head of a MIMIC-III discharge summary is the
administrative header block, where de-identification surrogates are densest. For
a substantial fraction of pairs the model is discriminating between two
administrative headers, which carry no information about whether the pairing is
genuine. A model given inputs that do not contain the discriminating signal
correctly learns to predict the base rate, and the base rate of a balanced
real-versus-mismatched task is exactly $0.5$.

\paragraph{The experiments that would separate them.} [TODO: (a) widen the
bottleneck from $d = 1$ to $d \in \{64, 256\}$ and rerun Stage 0, holding
everything else fixed; (b) hold the bottleneck fixed and replace head truncation
with truncation centred on a medication-bearing sentence, and rerun; (c) report
the fraction of truncated discharge-summary windows containing at least one
extracted medication mention, which bounds how often the input can possibly
carry signal.] Item (c) costs minutes and should be done first.

\section{Experimental Hygiene}
\label{sec:hygiene}

Three further defects concern how experiments were run rather than what was
built. They are reported because they bear on how much weight any reported
number can carry.

\subsection{Training on the evaluation set}

One reported configuration trained on the full $1{,}978$-pair pool, which
contains the $299$ pharmacist-annotated pairs used for evaluation. Results from
that configuration are not reported anywhere in this thesis. Any earlier project
artifact quoting them should be treated as superseded. The only reported
configuration is the one trained on the held-out complement.

\subsection{Seeds}

The ablation table reports three seeds with a standard deviation. In the code,
\texttt{SEEDS = [1, 2]}. The third row is a pasted literal from an earlier run at
a different code state. Further, \texttt{set\_seed()} is never called inside the
training loop, so each ablation variant inherits whatever RNG state the previous
variant left behind, making the variants sequentially dependent. The data split
uses \texttt{random\_state=42} in every case, so the reported $\pm$ reflects
initialisation noise only and contains no sampling variance whatever.

A separate problem compounds this: the full-model reference row in the ablation
table is itself a hardcoded literal. \textbf{The full model is not run by the
ablation script.} The table therefore compares three freshly trained variants
against a pasted reference, which means the differences it reports are not
differences measured under a common protocol. The ablation table should be
regenerated. Until it is, the only defensible reading of
Table~\ref{tab:acne} is directional.

\subsection{Row-level splits}

The reference pipeline split the candidate pool at the row level, as
Chapter~\ref{ch:candidates} documented: $1{,}978$ pairs drawn from $1{,}203$
admissions, roughly $1.64$ per admission, in a population with known repeat
admissions. Pairs sharing a nursing note, a discharge summary, or a patient
routinely land on both sides. The rebuild splits by patient with disjointness
asserted in code.

\subsection{The cumulative effect}

Taken together these mean that the ablation numbers in this chapter,
including Table~\ref{tab:acne}, should be read as descriptive rather than
inferential. They are reported because the qualitative pattern they show,
symmetric beating asymmetric, is informative regardless of the error bars. No
significance claim is made on any of them.

\section{A Diagnostic Protocol}
\label{sec:diagnostics}

The constructive contribution of this chapter is that each failure has a cheap,
mechanical detector, and that all of them belong in the training loop rather
than in a post-hoc audit. The protocol below is stated generally and each item
is annotated with the failure it catches.

\begin{enumerate}
\item \textbf{Print a module manifest at training time}: every named module,
      its parameter count, and a boolean indicating whether it appears in the
      autograd graph of the loss. \emph{Catches:} the PPC deletion, the
      untrained harm and explanation heads, and the $1.54$M-versus-$7.5$M
      parameter discrepancy (Section~\ref{sec:actual}).
\item \textbf{For every parameter reported as learned, log
      $\lVert\theta_{\text{final}} - \theta_{\text{init}}\rVert /
      \lVert\theta_{\text{init}}\rVert$ and a gradient-norm trace.}
      \emph{Catches:} the authority-conditioning artifact
      (Section~\ref{sec:c1}).
\item \textbf{For every auxiliary head, assert that its target is not derivable
      from its input features}, and verify by transfer rather than by held-out
      accuracy. \emph{Catches:} the rule-type leakage
      (Section~\ref{sec:c2}).
\item \textbf{Report a majority-class baseline alongside every $F_1$}, computed
      on the exact split being reported. \emph{Catches:} the uninterpretability
      of Table~\ref{tab:acne} under $84.9\%$-positive supervision.
\item \textbf{Assert group disjointness of splits in code}, at the coarsest
      grouping unit present in the data (here, patient). \emph{Catches:} the
      row-level split.
\item \textbf{Assert that the evaluation set is absent from the training set},
      by identifier, at load time. \emph{Catches:} the contaminated
      configuration.
\item \textbf{Call the seed setter inside the loop, not outside it}, and refuse
      to emit a standard deviation over fewer than three genuinely independent
      runs. \emph{Catches:} the seed reporting.
\item \textbf{Report the empirical range of every score the model emits.} A
      score spanning $0.398$ to $0.610$ across a full evaluation set is the
      range of an untrained sigmoid and should stop a results table.
      \emph{Catches:} the harm head (Section~\ref{sec:c3}).
\item \textbf{Run a transfer evaluation on any cohort at all before claiming
      generalisation.} \emph{Catches:} both the rule leakage and any surface
      memorisation.
\end{enumerate}

Every one of these is standard practice somewhere. None of them was in this
pipeline. The claim being made is not that they are novel; it is that their
absence is what allowed four separate unsupported claims to reach a submission,
and that a project of this shape should treat them as non-optional
instrumentation rather than as good hygiene.

\section{Proposed Figures}
\label{sec:arch-figures}

\begin{figure}[t]
\centering
\framebox[\linewidth]{\parbox[c][6cm][c]{0.92\linewidth}{\centering
\textit{[Figure asset 6.1: two-column architecture diagram. Left column,}
\textit{``as specified'': BioMedBERT $\rightarrow$ CDE $\rightarrow$ ATCG}
\textit{$\rightarrow$ PPC $\rightarrow$ GGHE $\rightarrow$ \{contradiction,}
\textit{rule, harm, explanation\} heads, with the 1.54M total. Right column,}
\textit{``as implemented'': the same skeleton with PPC struck through and}
\textit{replaced by a zero vector box, GGHE's bank annotated}
\textit{``randn(20,256)'', the harm and explanation heads greyed out and}
\textit{labelled ``no loss'', ACNE added at 5.9M, total 7.5M. Loss graph edges}
\textit{drawn only for the two heads that receive gradient.]}}}
\caption{Specification against implementation. The reader should notice that the
two heads carrying the project's severity and explainability claims are not
connected to the loss, and that the parameter budget is off by a factor of
roughly five in the opposite direction from what the writeups claim.}
\label{fig:arch}
\end{figure}

\begin{figure}[t]
\centering
\framebox[\linewidth]{\parbox[c][5cm][c]{0.92\linewidth}{\centering
\textit{[Figure asset 6.2: two panels. Left, a bar pair for the authority}
\textit{parameter: initialisation (0.5, 1.0) against final (0.509, 1.006),}
\textit{with the relative displacement annotated as 1.8\% and 0.6\%, plus a}
\textit{gradient-norm trace inset that is flat near zero. Right, rule-type F1}
\textit{across three cohorts: HF 0.9603, sepsis 0.0654, AKI 0.0595, with the}
\textit{45.2\% held-out rule accuracy marked as a horizontal reference.]}}}
\caption{The two falsifications in one view. Left: a parameter reported as
learned that did not move. Right: a head reported at $F_1 = 0.96$ that retains
nothing outside its training cohort. Both diagnostics are one line of code and
neither was run before the claims were made.}
\label{fig:falsification}
\end{figure}

\section{Limitations of This Chapter}
\label{sec:arch-limitations}

\paragraph{This is an audit, not a controlled study.} The falsifications are
established by reading code and by measurements already present in the run
artifacts. They establish that the claims were unsupported. They do not
establish what a correct implementation of any of the four components would do,
and this thesis makes no claim about that.

\paragraph{The Stage 0 mechanisms are hypotheses supported by code reading, not
by ablation.} Section~\ref{sec:stage0} is explicit about this and lists the
experiments that would separate them. The strongest evidence available is the
invariance across a $37.9\times$ capacity change, which rules out capacity but
does not adjudicate between the bottleneck and the truncation.

\paragraph{The transfer collapse has an alternative reading.} A drop from $0.96$
to $0.065$ is consistent with leakage plus genuine cohort shift, and the two are
not separated here. The candidate rules, the drug dictionary, and the class
dictionary are all heart-failure specific, so the sepsis and AKI candidate pools
differ in composition as well as in surface features. The leakage claim rests
primarily on the code, in which the rule identifier is an input feature; the
transfer number corroborates it rather than proving it alone.
[TODO: run the direct ablation, zero out \texttt{pf[i,1]} at inference and
re-measure in-domain rule $F_1$. If it collapses, the leakage claim is closed.]

\paragraph{The ablation numbers carry no valid error bars.}
Section~\ref{sec:hygiene}. They are reported directionally and no significance
claim is attached to any of them.

\paragraph{No comparison against a simple baseline architecture.} The project
never trained a plain BioMedBERT sentence-pair classifier on the same weak
labels and the same splits. That is the single most important missing
experiment in this chapter, because it would establish how much of the reported
behaviour is attributable to the custom components at all.
[TODO: train a vanilla cross-encoder baseline on the identical held-out split
and report it alongside every ACCORD row.]

\section{Summary}
\label{sec:arch-summary}

ACCORD was specified as a biomedical encoder plus four purpose-built components
totalling $1.54$M parameters. The model that produced every reported result
contains roughly $7.5$M new parameters, is missing one of the four components
entirely, and connects only two of its four output heads to the loss.

Three claimed contributions are falsified. The authority weights reported as
discovered were the initialisation values, displaced by under two percent, on a
branch that a subsequent LayerNorm largely erases; the project's own ablation
shows that removing the asymmetry improves both reported metrics. The
rule-type head's $F_1$ of $0.9603$ is a copy of an input feature, and it
collapses to $0.0654$ and $0.0595$ on two transfer cohorts and to $45.2\%$
accuracy on the held-out annotated set. Guideline-grounded harm estimation was
never trained: its guideline bank is twenty random vectors and its head is
absent from the loss, so its output range of $0.398$ to $0.610$ is the range of
an untrained sigmoid and the associated research question was not tested.

Self-supervised pretraining converged to the constant-prediction solution, at
$\ln 2$ with development accuracy exactly $0.5000$, identically under a
$37.9\times$ change in trainable capacity. The explanation originally given for
this, that clinical notes are too homogeneous, is withdrawn as untested. Two
mechanisms are visible in the code: a one-dimensional scoring bottleneck and
head truncation into the de-identification block.

Chapter~\ref{ch:eval} measures what this model does on the pharmacist-annotated
set, four independent ways, and finds that its scores contain no information
about the labels.
